% © 2026 Ing. David Jaroš — CC BY-NC-ND 4.0
%
% This work is licensed under a Creative Commons Attribution-NonCommercial-NoDerivatives
% 4.0 International License (CC BY-NC-ND 4.0).
%
% License History: Earlier drafts (up to v0.3) were released under CC BY 4.0.
% From v0.4 onward, all material is released under CC BY-NC-ND 4.0 to protect
% the integrity of the theoretical work during ongoing academic development.
%
% See LICENSE.md for full license text.

\ifdefined\INCLUDEMODE
\else
  \documentclass[12pt,a4paper]{article}
  \usepackage[a4paper, margin=2.5cm]{geometry}
  \usepackage{amsmath, amssymb, amsthm}
  \usepackage{mathtools}
  \usepackage{hyperref}
  \usepackage{xcolor}
  \usepackage{mdframed}
  \usepackage{booktabs}
  \usepackage{array}

  \newtheorem{theorem}{Theorem}[section]
  \newtheorem{lemma}[theorem]{Lemma}
  \newtheorem{proposition}[theorem]{Proposition}
  \newtheorem{corollary}[theorem]{Corollary}
  \theoremstyle{definition}
  \newtheorem{definition}[theorem]{Definition}
  \theoremstyle{remark}
  \newtheorem{remark}[theorem]{Remark}

  \newmdenv[backgroundcolor=green!15, linecolor=green!60!black, linewidth=1.5pt]{resultbox}
  \newmdenv[backgroundcolor=blue!8, linecolor=blue!50, linewidth=1.5pt]{insightbox}

  \title{\textbf{Matrix Representations of Discrete Symmetry Operators\\
                 in Biquaternion Field Theory}}
  \author{Ing.\ David Jaro\v{s}}
  \date{2026}

  \begin{document}
  \maketitle

  \begin{abstract}
  We provide explicit matrix realizations of the biquaternion algebra
  $\mathcal{B}=\mathbb{C}\otimes\mathbb{H}$ and its discrete symmetry operators
  $C$, $P$, $T_{\rm UBT}$, and $CPT$ in the standard $2\times2$ complex and
  $4\times4$ real representations.
  The purpose is to make UBT legible to particle physicists and mathematicians
  working in the matrix formalism.
  All objects defined abstractly in \texttt{discrete\_symmetries.tex} are
  re-expressed as concrete matrix operations.
  \end{abstract}
\fi

%──────────────────────────────────────────────────────────────────────────────
\section{The Standard $2\times2$ Complex Representation}
\label{sec:mat:2x2}
%──────────────────────────────────────────────────────────────────────────────

\subsection{Basis Elements}

Every biquaternion $b = b_0 + b_1\mathbf{I} + b_2\mathbf{J} + b_3\mathbf{K}$
with $b_a\in\mathbb{C}$ is mapped to a $2\times2$ complex matrix via the
\emph{standard representation} $\phi:\mathcal{B}\to M_2(\mathbb{C})$:
\begin{equation}
  \phi(1) = \begin{pmatrix} 1 & 0 \\ 0 & 1 \end{pmatrix},
  \quad
  \phi(\mathbf{I}) = \begin{pmatrix} i & 0 \\ 0 & -i \end{pmatrix},
  \quad
  \phi(\mathbf{J}) = \begin{pmatrix} 0 & 1 \\ -1 & 0 \end{pmatrix},
  \quad
  \phi(\mathbf{K}) = \begin{pmatrix} 0 & i \\ i & 0 \end{pmatrix}.
  \label{eq:mat:basis_2x2}
\end{equation}
Note: $\phi(\mathbf{I}) = i\sigma_3$, $\phi(\mathbf{J}) = i\sigma_2$,
$\phi(\mathbf{K}) = i\sigma_1$, where $\sigma_k$ are the standard Pauli matrices.

A general biquaternion maps to:
\begin{equation}
  \phi(b) = \begin{pmatrix}
    b_0 + ib_1 & b_2 + ib_3 \\
    -b_2 + ib_3 & b_0 - ib_1
  \end{pmatrix} \in M_2(\mathbb{C}).
  \label{eq:mat:general_2x2}
\end{equation}

\begin{lemma}[Faithfulness of $\phi$]
$\phi$ is a faithful $\mathbb{C}$-algebra homomorphism, i.e., injective and
compatible with multiplication and addition.
\end{lemma}
\begin{proof}
Quaternion multiplication rules are verified by direct computation:
$\phi(\mathbf{I})\phi(\mathbf{J}) = (i\sigma_3)(i\sigma_2) = -\sigma_3\sigma_2
= -\begin{pmatrix}1&0\\0&-1\end{pmatrix}\begin{pmatrix}0&-i\\i&0\end{pmatrix}
= \begin{pmatrix}0&i\\i&0\end{pmatrix} = \phi(\mathbf{K})$. \checkmark
Injectivity: the image is the $\mathbb{C}$-span of
$\{I_2, i\sigma_3, i\sigma_2, i\sigma_1\}$, which are linearly independent.
\end{proof}

\subsection{Involutions as Matrix Operations}

The algebraic involutions of Section~2 in \texttt{discrete\_symmetries.tex}
translate to matrix operations as follows:

\paragraph{$P_1$ (complex conjugation):}
\begin{equation}
  \phi(P_1(b)) = \overline{\phi(b)} = \begin{pmatrix}
    \bar b_0 - i\bar b_1 & \bar b_2 - i\bar b_3 \\
    -\bar b_2 - i\bar b_3 & \bar b_0 + i\bar b_1
  \end{pmatrix}.
  \label{eq:mat:P1_matrix}
\end{equation}
At the matrix level: $P_1$ acts as entry-wise complex conjugation of the
matrix $\phi(b)$, i.e., $\phi(P_1(b)) = \overline{\phi(b)}$.

\paragraph{$P_2$ (quaternion conjugation):}

Direct computation from \eqref{eq:mat:general_2x2} gives:
\begin{equation}
  \phi(P_2(b))
  = \begin{pmatrix} b_0 - ib_1 & -b_2 - ib_3 \\ b_2 - ib_3 & b_0 + ib_1 \end{pmatrix}.
  \label{eq:mat:P2_matrix}
\end{equation}
Comparing with $\phi(b) = \begin{pmatrix}a & b \\ c & d\end{pmatrix}$
(where $a = b_0+ib_1$, $b = b_2+ib_3$, $c = -b_2+ib_3$, $d = b_0-ib_1$),
this is the \emph{reverse} operation on the matrix:
\begin{equation}
  \phi(P_2(b)) = \phi(b)^{\rm rev},
  \qquad
  M^{\rm rev} := \begin{pmatrix} M_{22} & -M_{12} \\ -M_{21} & M_{11} \end{pmatrix}.
  \label{eq:mat:P2_rev}
\end{equation}
For unit biquaternions ($\det\phi(b) = 1$), $M^{\rm rev} = M^{-1}$.
For real quaternions with real $b_a\in\mathbb{R}$ (so $\phi(b)\in SU(2)$,
$\det\phi(b)=1$, and $\phi(b)^\dagger = \phi(b)^{-1}$), $P_2 \leftrightarrow (\ )^\dagger$
holds.  For general complex biquaternions, one must use
$P_2 \leftrightarrow (\ )^{\rm rev}$ \eqref{eq:mat:P2_rev}, not $(\ )^\dagger$.

\begin{remark}[Why $P_2\neq(\ )^\dagger$ for complex biquaternions]
$\phi(b)^\dagger = \phi(P_{12}(b))$ (Hermitian adjoint corresponds to $P_{12}$,
not $P_2$, as shown below).
Setting $\phi(P_2(b)) = \phi(b)^\dagger$ would require $b_a = \bar b_a$, i.e.,
real coefficients.  The identification $P_2\leftrightarrow(\ )^\dagger$ is
therefore valid only for the $\mathbb{R}$-subalgebra $\mathbb{H}\subset\mathcal{B}$.
\end{remark}

\paragraph{$P_{12}$ (Hermitian conjugation = $P_1\circ P_2$):}

Direct computation:
\begin{equation}
  \phi(P_{12}(b)) = \phi(\bar b_0 - \bar b_1\mathbf{I}
                     - \bar b_2\mathbf{J} - \bar b_3\mathbf{K})
  = \begin{pmatrix}
    \bar b_0 - i\bar b_1 & -\bar b_2 - i\bar b_3 \\
    \bar b_2 - i\bar b_3 & \bar b_0 + i\bar b_1
  \end{pmatrix}
  = \phi(b)^\dagger.
  \label{eq:mat:P12_matrix}
\end{equation}
Hence the \emph{Hermitian adjoint} $(\ )^\dagger$ corresponds to the full
Hermitian conjugation $P_{12} = P_1\circ P_2$, not to $P_2$ alone.
This is the corrected identification: for \emph{general} biquaternions,
$(\ )^\dagger \leftrightarrow P_{12}$, and
$(\ )^{\rm rev} \leftrightarrow P_2$.

%──────────────────────────────────────────────────────────────────────────────
\section{Symmetry Operators in the $2\times2$ Representation}
\label{sec:mat:sym_2x2}
%──────────────────────────────────────────────────────────────────────────────

In the $2\times2$ representation, the field $\Theta(q,\tau)$ is identified with
a spacetime- and $\tau$-dependent matrix
$\mathbf{\Theta}(q,\tau) = \phi(\Theta(q,\tau)) \in M_2(\mathbb{C})$.

\subsection{Charge Conjugation $C$}

\begin{equation}
  C[\mathbf{\Theta}](q,\tau) = \overline{\mathbf{\Theta}(q,\tau^*)}.
  \label{eq:mat:C_2x2}
\end{equation}
Action: complex-conjugate all matrix entries AND replace $\tau\to\tau^*$.

For gauge fields in the matrix representation:
$C: \mathbf{A}_\mu \mapsto -\mathbf{A}_\mu^*$ (charge reversal).

\subsection{Spatial Parity $P$}

\begin{equation}
  P[\mathbf{\Theta}](q,\tau) = \mathbf{\Theta}(\bar q,\tau)^{\rm rev},
  \label{eq:mat:P_2x2}
\end{equation}
i.e., the \emph{reverse} of the matrix field (Definition~\eqref{eq:mat:P2_rev})
evaluated at the reflected spacetime point $\bar q = (q^0,-\mathbf{q})$.
For real quaternions (or unit biquaternions), this equals $\mathbf{\Theta}(\bar q,\tau)^{-1}$.

\subsection{Time Reversal $T_{\rm UBT}$}

\begin{equation}
  T_{\rm UBT}[\mathbf{\Theta}](q,\tau) = \mathbf{\Theta}(q,-\tau).
  \label{eq:mat:T_2x2}
\end{equation}
No matrix operation on $\mathbf{\Theta}$; only $\tau\to-\tau$.

\subsection{$CPT$ Operator}

\begin{equation}
  CPT[\mathbf{\Theta}](q,\tau)
  = (C\circ P\circ T_{\rm UBT})[\mathbf{\Theta}](q,\tau)
  = \overline{\mathbf{\Theta}(\bar q,-\tau)^{\rm rev}}
  = \mathbf{\Theta}(\bar q,-\tau)^\dagger.
  \label{eq:mat:CPT_2x2}
\end{equation}
The last equality uses $\overline{M^{\rm rev}} = M^\dagger$ for any $2\times2$
complex matrix (direct verification from definitions).
In the $2\times2$ representation, $CPT$ acts as the \emph{Hermitian adjoint}
$(\ )^\dagger$ combined with $(q,\tau)\to(\bar q,-\tau)$.

\begin{resultbox}
\textbf{Summary: symmetry operators in $2\times2$ representation.}
\begin{center}
\begin{tabular}{@{}lll@{}}
\toprule
Operator & Matrix operation on $\mathbf{\Theta}$ & Argument action \\
\midrule
$C$ & $\mathbf{\Theta} \mapsto \overline{\mathbf{\Theta}}$ (entry-wise conj.) & $\tau\to\tau^*$ \\
$P$ & $\mathbf{\Theta} \mapsto \mathbf{\Theta}^{\rm rev}$ (reverse, eq.~\eqref{eq:mat:P2_rev}) & $q\to\bar q$ \\
$T_{\rm UBT}$ & $\mathbf{\Theta} \mapsto \mathbf{\Theta}$ (identity) & $\tau\to -\tau$ \\
$CPT$ & $\mathbf{\Theta} \mapsto \mathbf{\Theta}^\dagger$ (Hermitian adj.) & $(q,\tau)\to(\bar q,-\tau)$ \\
\bottomrule
\end{tabular}
\end{center}
Note: $P$ maps $\mathbf{\Theta}\to\mathbf{\Theta}^{\rm rev}$ (not $(\ )^\dagger$) for general
biquaternions; for real quaternions $(\det\mathbf{\Theta}=1)$ the two coincide.
The $CPT$ operator corresponds to the Hermitian adjoint $(\ )^\dagger$ because
$CPT = P_{12}$ at the algebra level ($P_1\circ P_2$, see \eqref{eq:mat:P12_matrix}).
\end{resultbox}

%──────────────────────────────────────────────────────────────────────────────
\section{The $4\times4$ Real Representation}
\label{sec:mat:4x4}
%──────────────────────────────────────────────────────────────────────────────

\subsection{Basis}

Treating $b_a = b_a^r + ib_a^i$ ($b_a^r, b_a^i\in\mathbb{R}$),
a biquaternion has 8 real components
$(b_0^r, b_0^i, b_1^r, b_1^i, b_2^r, b_2^i, b_3^r, b_3^i)$.
The left-regular representation maps $\mathcal{B}\to M_8(\mathbb{R})$.

For the \emph{quaternion-only} part (four real components), the
$4\times4$ real representation of $\mathbb{H}$ via left-multiplication
by unit quaternions is:
\begin{equation}
  \phi_{4}(1) = I_4,
  \quad
  \phi_{4}(\mathbf{I}) = \begin{pmatrix}
    0 & -1 & 0 & 0 \\
    1 &  0 & 0 & 0 \\
    0 &  0 & 0 & -1 \\
    0 &  0 & 1 &  0
  \end{pmatrix},
  \quad
  \phi_{4}(\mathbf{J}) = \begin{pmatrix}
    0 & 0 & -1 & 0 \\
    0 & 0 &  0 & 1 \\
    1 & 0 &  0 & 0 \\
    0 & -1 & 0 & 0
  \end{pmatrix},
  \quad
  \phi_{4}(\mathbf{K}) = \begin{pmatrix}
    0 & 0 & 0 & -1 \\
    0 & 0 & -1 & 0 \\
    0 & 1 & 0 & 0 \\
    1 & 0 & 0 & 0
  \end{pmatrix}.
  \label{eq:mat:4x4_basis}
\end{equation}
These are block-antisymmetric orthogonal matrices satisfying the quaternion
algebra.

\subsection{Involutions in $4\times4$ Representation}

\paragraph{$P_2$ (quaternion conjugation):}
In the $4\times4$ representation, quaternion conjugation sends
$\mathbf{e}_i\to -\mathbf{e}_i$ ($i=1,2,3$):
\begin{equation}
  \phi_4(P_2(b)) = \mathrm{diag}(+1,-1,-1,-1)\,\phi_4(b)\,\mathrm{diag}(+1,-1,-1,-1).
  \label{eq:mat:P2_4x4}
\end{equation}
This is conjugation by the ``metric signature'' matrix
$\eta = \mathrm{diag}(1,-1,-1,-1)$:
$\phi_4(P_2(b)) = \eta\,\phi_4(b)\,\eta$ (since $\phi_4$ is the left-regular
representation and $P_2$ reverses the sign of the vector part).

\paragraph{Spatial parity $P$ in $4\times4$ matrix language:}
The field value transforms as
$P[\Theta](q,\tau) = P_2(\Theta(\bar q,\tau))$, which in the $4\times4$
representation reads:
\begin{equation}
  P[\mathbf{\Theta}_{4}](q,\tau)
  = \eta\,\mathbf{\Theta}_{4}(\bar q, \tau)\,\eta.
  \label{eq:mat:P_4x4}
\end{equation}
This has the structure of a Lorentz reflection on the internal space.

%──────────────────────────────────────────────────────────────────────────────
\section{Worked Example: $CP$ Transformation}
\label{sec:mat:example_CP}
%──────────────────────────────────────────────────────────────────────────────

Let $\Theta$ be a biquaternion field with purely complex-scalar part:
$\Theta = (a + ib)\mathbf{1}$, $a,b\in\mathbb{R}$.  In the $2\times2$
representation:
\begin{equation}
  \mathbf{\Theta} = (a+ib)I_2.
\end{equation}

\paragraph{Under $C$:}
$C[\mathbf{\Theta}] = \overline{(a+ib)}I_2 = (a-ib)I_2$. \checkmark
(Complex conjugation flips the phase.)

\paragraph{Under $P$:}
$P[\mathbf{\Theta}](\bar q) = \mathbf{\Theta}(\bar q)^{\rm rev}$.
For $\mathbf{\Theta} = (a+ib)I_2$: $M^{\rm rev} = \begin{pmatrix}a+ib & 0\\0 & a+ib\end{pmatrix} = (a+ib)I_2$.
So $P[\mathbf{\Theta}] = (a+ib)I_2$ at $\bar q$.  A pure-scalar field is $P$-invariant. \checkmark

\paragraph{Under $CP$:}
$CP[\mathbf{\Theta}](q,\tau^*) = \mathbf{\Theta}(\bar q,\tau^*)^\dagger$.
For $\mathbf{\Theta} = (a+ib)I_2$: $((a+ib)I_2)^\dagger = (a-ib)I_2$.
If $b=0$ (real scalar), $CP[\mathbf{\Theta}] = (a+ib)I_2$: $CP$-invariant. \checkmark

\begin{table}[h]
\centering
\caption{Involution action: algebra vs.\ $2\times2$ matrix representation.
  Notation: $M = \phi(b)$; $b = b_0\mathbf{1} + b_1\mathbf{I} + b_2\mathbf{J} + b_3\mathbf{K}$,
  $b_a\in\mathbb{C}$ general.}
\label{tab:mat:check}
\begin{tabular}{@{}lll@{}}
\toprule
Involution & Biquaternion result & $2\times2$ matrix identification \\
\midrule
$P_1(b)$ & $\bar b_0 + \bar b_1\mathbf{I} + \bar b_2\mathbf{J} + \bar b_3\mathbf{K}$
  & $\phi(P_1(b)) = \overline{M}$ \quad (exact for all $b$) \\
$P_2(b)$ & $b_0 - b_1\mathbf{I} - b_2\mathbf{J} - b_3\mathbf{K}$
  & $\phi(P_2(b)) = M^{\rm rev}$ \quad (exact for all $b$)\\
  & & \small{where $M^{\rm rev} = \begin{pmatrix}M_{22} & -M_{12}\\-M_{21} & M_{11}\end{pmatrix}$}\\
  & & \small{Special case: $b_a\in\mathbb{R}$ $\Rightarrow$ $M^{\rm rev} = M^\dagger$} \\
$P_{12}(b)$ & $\bar b_0 - \bar b_1\mathbf{I} - \bar b_2\mathbf{J} - \bar b_3\mathbf{K}$
  & $\phi(P_{12}(b)) = M^\dagger$ \quad (exact for all $b$) \\
\bottomrule
\end{tabular}
\end{table}

\begin{insightbox}
\textbf{Canonical matrix identification.}
For general biquaternions with complex coefficients $b_a\in\mathbb{C}$:
\begin{center}
\begin{tabular}{ll}
$P_1 \leftrightarrow \overline{M}$ & (entry-wise conjugation) \\
$P_2 \leftrightarrow M^{\rm rev}$ & (reverse/adjugate-like operation) \\
$P_{12} \leftrightarrow M^\dagger$ & (Hermitian adjoint) \\
\end{tabular}
\end{center}
The identification $P_2\leftrightarrow(\ )^\dagger$ and
$P_{12}\leftrightarrow(\ )^T$ that appear in some quaternion-algebra texts
hold only for \textbf{real quaternions} ($b_a\in\mathbb{R}$).
For full biquaternions, one must use the assignments above.
\end{insightbox}

%──────────────────────────────────────────────────────────────────────────────
\section{Chiral Projectors in Matrix Form}
\label{sec:mat:chiral}
%──────────────────────────────────────────────────────────────────────────────

The left/right chiral projectors
$\mathsf{P}_\pm = \frac{1}{2}(\mathbf{1}\pm P_2)$
(Section~2.1 of \texttt{chirality\_and\_parity\_breaking.tex})
act in the $2\times2$ representation as:
\begin{equation}
  \phi(\mathsf{P}_+b) = \frac{1}{2}\left(M + M^{\rm rev}\right)
  = \begin{pmatrix} b_0 & 0 \\ 0 & b_0 \end{pmatrix} = b_0 I_2
  \quad\text{(scalar part)},
  \label{eq:mat:P_plus}
\end{equation}
\begin{equation}
  \phi(\mathsf{P}_-b) = \frac{1}{2}\left(M - M^{\rm rev}\right)
  = \begin{pmatrix} ib_1 & b_2+ib_3 \\ -b_2+ib_3 & -ib_1 \end{pmatrix}
  = ib_1\sigma_3 + ib_3\sigma_2 + ib_2\sigma_1
  \quad\text{(vector part)}.
  \label{eq:mat:P_minus}
\end{equation}
The left-handed sector $\mathsf{P}_-$ selects the \emph{traceless anti-Hermitian}
part of the matrix (for $b_a\in\mathbb{R}$); for complex $b_a$ it selects
the off-scalar (non-central) degrees of freedom.
Traceless anti-Hermitian $2\times2$ matrices form the Lie algebra
$\mathfrak{su}(2)$, consistent with the identification of $\Theta_L$
with the $SU(2)_L$ fundamental representation.

\ifdefined\INCLUDEMODE
\else
  \end{document}
\fi
